\newpage
\section*{About the Author}
\addcontentsline{toc}{chapter}{About the Author}
Jueon Kim and Yeoeun Kim are students who have lived in Korea, China, and the 
United States, where they experienced firsthand how language and culture shape 
learning. Adapting to new environments and study methods brought many challenges, 
but guided by their parents' teaching to \inlinequote{learn well and share generously,} 
they embraced service as a core value.

Through their nonprofit organization GYCO, they bring music to seniors, the 
homeless, and cancer patients each month, and send books to refugees and students 
in Africa to support education. After performances, they talk with many seniors 
to help ease loneliness.

Their commitment to service has inspired many peers, and this math book reflects 
their own journey as Korean immigrants adapting to new languages and education 
systems. They hope it helps students overcome language barriers, see math as 
approachable, and use it as a tool for confidence and opportunity.

\subsection*{저자 소개}
김주언, 김여은 두 학생은 한국, 중국, 미국에서 생활하며 언어와 문화가 학습에 어떤 
영향을 주는지 직접 경험한 학생입니다. 이 학생들은 변화된 환경에 적응해야 했고 새로운 
문화와 학습 방법을 받아들이는 과정 속에서 어려움도 많았습니다.

부모님께 배운 \inlinequote{잘 배워 잘 나누는 사람이 되라}는 가르침은 이들의 삶의 중요한 
가치가 되었습니다. 비영리단체 GYCO에서 사회의 소외계층에게 사랑과 희망을 전하는 다양한 
활동을 현재까지 꾸준하게 하고 있습니다.

음악 연주를 직접 보러 올 수 없으신 사회의 연약한 계층(시니어, 노숙자, 암환자)에게 한 
달에 한 번씩 연주로 찾아가고, 교육이 필요한 난민과 아프리카의 학생들에게 책을 보내 
교육받을 수 있게 돕고 있습니다. 많은 시니어들에게 찾아가 연주가 끝나고 함께 이야기를 
나누며 그들의 외로움을 덜어드렸고, 노숙자분들께 필요한 물품을 준비해서 나눠드리는 활동도 
지속적으로 해오고 있습니다. 함께 봉사하는 학생들에게도 이들의 성실함과 열정과 나눔은 늘 
도전이 되고 있습니다.

이 수학책은 한국 이민자로서 새로운 언어와 교육 제도에 적응하면서 느낀 두 학생의 어려움을 
바탕으로, 같은 길을 걷는 학생들이 수학을 더 쉽고 친근하게 느끼도록 돕기 위해 시작되었습니다.

언어의 장벽으로 힘들어 하는 학생들이 수학을 어려움으로 느끼지 않고, 변화된 상황에서 뛰어난 
수학 실력이 그들에게 자신감을 갖고 기회를 넓혀주는 도구가 되기를 희망하며 책을 만든 
귀한 학생들입니다.